电力电子技术作为陪伴我大半个本科生活的一门学科，可以说就像《一块红布》中描写的那样差点把我闷死。

大二的时候，我直接零基础参加大创比赛，和队友靠着拼凑论文尸块得出的算法，以及使用半生不熟的 Simulink 技术做的仿真，让我们的 DAB 项目拿了校内二等奖，后又拿了上海市交科赛的三等奖，最后无疾而终了（笑）。

大三的时候，我们正式学习了电力电子技术这门课，用的小众冷门变态教材（国内用王兆安老先生的比较多，我们用的不是这本），讲课也讲得毫无逻辑，因为课比较多，能做完作业就差不多得了，所以课后没额外花时间，学也学得一坨，虽然最后拿了优，但是是纯靠记的。一年前竞赛中的状态分析、纹波计算依旧不会，或者说换一个拓扑就不会了，可见学的确实是一坨。

虽然我对电力电子有点感兴趣，但是学了一个学期之后兴趣直接没了啊嗯。可以说电力电子究竟是个什么样子，我搞了一个学期，也并没有完全搞清楚。因此在读研之前，我决定重新学习电力电子技术这门课，彻彻底底地把原来学的那一坨给弄清楚。

“大学大学，大不了自学。”大一分专业之前的班主任如是说。本来没放在心上，在被专业课拷打了两年之后觉得说的太好了，我完全同意。于是我就开始在 B 站上搜教程，简单听了几节教程，感觉还是王兆安老先生和上科大傅旻帆老师的最好，其中又以傅旻帆老师的最甚。

为什么呢？因为我最希望的就是系统性掌握分析方法，多来点干货，而不是一上来就讲器件，讲了我又记不住，讲来干嘛，，，傅老师的课程就比较对我胃口，上来就讲方法，从稳态量到纹波分析，一环扣一环，我现在才明白原来分析纹波不是靠玄学，，，当然后面也讲了器件，但是是基于一定的逻辑，即将想象中的开关转化为实际的器件，并针对器件的参数做了全面的建模，并应用前面所学知识进行分析。

总之，在听了傅老师的课之后，我仿佛是醍醐灌顶，心中又燃起了学习电力电子的斗志，，，虽然之后主要是做电力系统方向，但是作为电气基础课，电力电子得学。当然电力系统分析肯定也是一个重量级，也得学，所以就准备和电力电子一起学了。学完这俩，还得补一下铁路牵引供电系统，如果有必要还要学一些现代控制理论，今年还得做毕设，真是充实的一年啊（笑）。

哦对了， GPT 大人太厉害了，把我笔记的图直接画到 \LaTeX{} 里面，我都不敢想纯手画得有多折磨，，，

碎碎念就到这里了，时间紧任务重，大大方方学八分钱，大踏步走进研究生生活！
